\subsection*{Partie A}
\begin{enumerate}
	\item La longueur totale d'un bonbon est de \np[mm]{15}.

	À chaque extrémité, il y a une demi-boule de rayon \np[mm]{2.5}.

	Donc la partie cylindrique a une hauteur égale à {$15-2\times\np{2.5}=\np[mm]{10}$.}

	\item \begin{enumerate}
		\item La partie cylindrique a pour base un disque de rayon \np[mm]{2.5},

		son aire est égale à {$\pi \times\np{2.5}^2\approx\np[mm^2]{19.63}$}.

		La hauteur du cylindre est de \np[mm]{10},

		donc son volume est égal à {$10\times(\pi\times\np{2.5}^2)\approx \np[mm^3]{196}$.}

		\item Les deux demi-boules qui constituent les extrémités ont, ensemble, le volume d'une boule de rayon \np{2.5}, soit \quad {$\dfrac43 \times \pi\times \np{2.5}^3\approx \np[mm^3]{65.4}$}.

		Léa a donc raison, puisque la somme des volumes de la partie cylindrique et des deux extrémités est égale à {$\np{196.3}+\np{65.4}=\np{261.7}$} qui s'arrondit à \np[mm^3]{262}.

		\textit{Variante:} Le volume des deux demi-boules est environ \np[mm^3]{65}, le volume total d'un bonbon est de \np[mm^3]{261}, donc Léa a raison.
	\end{enumerate}

	\item Un bonbon demande un volume d'environ \np[mm^3]{262},

	donc pour produire \np{300000} bonbons, il faut {$\np{300000}\times262=\np[mm^3]{78600000}$} de mélange.

	Or \quad {$\np[litre]{1} = \np[dm^3]{1} = \np[cm^3]{1000} = \np[mm^3]{1000000}$}.

	Donc \quad  {$\np[mm^3]{78600000}=\np[litres]{78.6}$} :

	83~litres seront donc plus que suffisants pour produire \np{300000}~bonbons.
\end{enumerate}

\subsection*{Partie B}

	Avec le format A, Léa paiera {$2\times\np{7.90}=\np[\euro{}]{15.80}$} pour 1 kilo de bonbons.

	Avec le format B, il lui faudra acheter 4 sachets, les 3 premiers à \np[\euro{}]{4.30}, le quatrième à \np[\euro{}]{2.15}.

	Cela lui coûtera au total \quad {$3\times\np{4.30}+\np{2.15}=\np[\euro{}]{15.05}$}.

	Léa a donc intérêt à profiter de la promotion sur les sachets B.

\needspace{10\baselineskip}
